\chapter{The Conditional Information-Spinor Critical-Axis Theorem}
\label{ch:conditional}

\section{Abstract cancellation representation}
We now state the exact conditional result. Let $\rho=\sigma+it$ be a nontrivial zero of $\xi$. Suppose there exists a canonical normalized representation of its zero condition by two complementary branches. The word \emph{canonical} means that the representation is derived from zeta data and does not contain parameters fitted separately to each zero.

\begin{assumption}[Complementary zero-state representation]
\label{ass:zero-state}
For each nontrivial zero $\rho$ there exist normalized branch vectors $\ket{0_\rho}$ and $\ket{1_\rho}$, a nonzero scalar normalization $N(\rho)$, and a readout functional $L_\rho$ symmetric under branch exchange such that
\begin{equation}
\xi(\rho)=N(\rho)L_\rho\left[
\sqrt{\sigma}\ket{0_\rho}
+e^{i\phi(\rho)}\sqrt{1-\sigma}\ket{1_\rho}
\right],
\label{eq:zero-rep}
\end{equation}
with the following properties:
\begin{enumerate}[label=(A\arabic*)]
\item $\sigma=\Ree\rho$;
\item $J:\rho\mapsto1-\overline\rho$ exchanges the two branches;
\item symmetry of $L_\rho$ reduces zerohood to cancellation of equal readout amplitudes;
\item the relative zero phase is $\phi(\rho)=\pi\pmod{2\pi}$;
\item $N(\rho)\neq0$ and neither branch readout vanishes separately at $\rho$.
\end{enumerate}
\end{assumption}

Condition (A5) blocks a tautological factorization in which each branch already contains the factor $\xi(\rho)=0$.

\begin{theorem}[Conditional critical-axis selection]
\label{thm:conditional}
Under Assumption \ref{ass:zero-state}, every nontrivial zero $\rho$ satisfies
\begin{equation}
\Ree\rho=\frac12.
\end{equation}
\end{theorem}

\begin{proof}
By (A3)--(A5), the zero condition is a genuine cancellation between two nonzero complementary amplitudes. By (A4), their relative phase is $\pi$. Symmetric readout then reduces the condition to the amplitude of Theorem \ref{thm:cancellation}:
\begin{equation}
\sqrt{\sigma}-\sqrt{1-\sigma}=0.
\end{equation}
Hence $\sigma=1/2$. By (A1), $\sigma=\Ree\rho$, giving the result.
\end{proof}

\begin{corollary}
If Assumption \ref{ass:zero-state} is established for every nontrivial zero, then the Riemann hypothesis follows.
\end{corollary}

This corollary is logically correct but does not lessen the difficulty: proving the assumption is the substantive open task.

\section{Equivalent fixed-point formulation}
The same result can be expressed without explicit amplitudes.

\begin{assumption}[Local self-duality]
Each nontrivial zero is individually invariant under the anti-linear zeta involution:
\begin{equation}
J(\rho)=\rho.
\end{equation}
\end{assumption}

Then the fixed-set theorem of Chapter \ref{ch:zeta} immediately gives $\Ree\rho=1/2$. However, this assumption is nearly equivalent to RH itself. The amplitude formulation is more useful only if the local self-duality can be derived from independent normalization, unitarity, or positivity principles.

\section{Information and holonomy version}
A more geometric set of sufficient conditions is:
\begin{enumerate}
\item zero states are normalized points of a two-state projective space;
\item complement acts as the exchange of populations;
\item analytic continuation around the relevant cycle has holonomy $-1$;
\item the zero readout is orthogonality to the symmetric state;
\item the representation is faithful and nondegenerate.
\end{enumerate}
The holonomy condition yields $\sigma=1/2$ modulo one, while the probability domain selects the interior value $1/2$. Shannon entropy at that point is then $\ln2$.

\begin{proposition}[Triple coincidence under the model]
Under Assumption \ref{ass:sigma}, the following conditions are equivalent within the two-state representation:
\begin{align}
\sigma&=\frac12,\\
\Hbin(\sigma)&=\ln2,\\
C(\sigma)&=\sigma,\\
n_z&=0,\\
\mathcal U_B(\sigma)&=-1,\\
\mathcal A(\sigma,\pi)&=0.
\end{align}
\end{proposition}

\begin{proof}
The equivalence of the first three statements follows from Chapter \ref{ch:shannon}; the Bloch condition from $n_z=2\sigma-1$; the holonomy from Chapter \ref{ch:bloch}; and the cancellation condition from Theorem \ref{thm:cancellation}.
\end{proof}

This proposition is a genuine unification inside the model. It says that maximum binary distinction, complement fixedness, equatorial geometry, spinor sign, and destructive interference are different coordinates of the same balanced state.

\section{Why off-axis quartets remain possible without the assumption}
Suppose $\rho=1/2+\delta+it$ with $\delta\neq0$ is a hypothetical zero. The zeta symmetries generate the partner $J(\rho)=1/2-\delta+it$. A global two-state construction could assign those two distinct points as complementary states rather than forcing either one to be self-dual. In that case the pair as a whole is balanced even though each member is not.

Therefore a proof must exclude \emph{pairwise balance without pointwise balance}. This is the exact obstruction. Possible exclusion mechanisms include:
\begin{enumerate}
\item a self-adjoint spectral operator whose eigenparameter is real;
\item a positive norm that would be violated by paired nonreal eigenvalues;
\item a local unitarity law requiring $|\chi(s)|=1$ at every zero;
\item a canonical readout for which branch norms cannot be externally rescaled;
\item a total-positivity theorem for the Xi kernel.
\end{enumerate}

\section{A no-circularity criterion}
Any future proof claiming to establish Assumption \ref{ass:zero-state} should pass the following test:
\begin{quote}
Could the proposed branch states and normalization be defined before knowing that $\Ree\rho=1/2$, and do they remain nonsingular for a hypothetical off-axis zero?
\end{quote}
If the construction explicitly inserts $\sigma=1/2$, uses a list of critical-line zeros as input, or divides by a factor vanishing at the target zero, it is circular.
